
Models of neutrino interactions are improving rapidly and experiments must make a choice of models to use for efficiency and background estimation at the time of an analysis.  
For this MicroBooNE measurement, GENIE v2 is used.  
However, data are compared with GENIE v3.0.6 with tune G18\_10a\_02\_11a (labelled as ``GENIE v3''), NuWro 19.02.1~\cite{Golan:2012wx}, NEUT 5.4.0.1~\cite{Hayato:2009zz}, and GiBUU 2019~\cite{Buss:2011mx}.  
As described in Sec.~\ref{sec:xsect_def}, each calculation was folded with our smearing matrix before comparing with data.

GENIE, NuWro, and NEUT are developed by separate groups who often implement the same models in different ways.
In this case, each code uses the local Fermi gas (LFG) momentum distribution and a binding energy correction which can be constant (GENIE and NEUT) or derived from a potential (NuWro) for sampling the struck nucleon properties.  
Each code uses a similar CCQE model, but apply different empirical RPA corrections.  
GENIE also applies a Coulomb correction~\cite{Nieves:2004wx} on the outgoing muon.
All modern codes contain a multinucleon mechanism.  
GENIE v3, NuWro, and NEUT all use the Valencia $2p2h$ model~\cite{Gran:2013kda}.
The new models are built to agree with \mb data~\cite{miniboone-ccqe}, which uses the same beamline as \mub but a hydrocarbon target.

Both NEUT and GENIE use the Kuzmin-Lubushkin-Naumov~\cite{Kuzmin:2003ji} and Berger-Sehgal\cite{Berger:2007rqRES} model for  resonance and the Berger-Sehgal model for coherent~\cite{Berger:2008xsCOH} scattering.
Implementation differs in the cutoff in total hadronic energy in the nucleon rest frame $W$ between resonance and DIS models, e.g. fixed cutoff value of 1.9 GeV (GENIE) and 2.0 GeV (NEUT).  NuWro is a little different, using the Adler-Rarita-Schwinger model~\cite{Graczyk:2007xk} for the $\Delta$(1232) resonance and a smooth transition to DIS at 1.6 GeV.
All use form factors fit to data.  
In addition, each code uses a similar extrapolation of DIS processes to $\pi N$ threshold~\cite{Yang:1999xg,Bodek:2002ps} to model non-resonant processes.
Hadrons produced at the neutrino interaction vertex undergo FSI which are governed by models based on the impulse approximation with nuclear medium corrections. Here, FSI includes both pion and nucleon interactions with most generators using separate models to simulate the two types of interaction.
NEUT and NuWro use the Salcedo-Oset model~\cite{Salcedo:1987md} for pions which includes nuclear medium effects.  
GENIE uses a more empirical model~\cite{Katori:2013eoa} for pions which has no nuclear medium corrections.  
NuWro and GENIE use nuclear medium corrections for nucleon FSI~\cite{Pandharipande:1992zz} while the NEUT code is strictly impulse approximation for nucleon FSI.

GiBUU~\cite{Buss:2011mx} was developed from a first principles approach as opposed to the greater emphasis on empirical models with other generators.
Although the interactions covered are very similar to the generators described above, the details are different as all of the components were developed as a coherent package.
Instead of a semi-classical approximation for propagating hadrons as is used in the models described above, a transport equation is solved~\cite{Buss:2011mx}.
Nucleon momenta come from an LFG distribution, then hadrons propagate through the residual nucleus in a nuclear potential that is consistent between initial and final state.
The CCQE interaction is handled with a spectral function that has separate momentum and energy dependence.
All resonances propagate as particles and interact according to best available models.   
Version 2019 is used here; a recent publication~\cite{Mosel:2017ssx} shows comparisons of this version with the T2K $CC0\pi$ data~\cite{Abe:2018pwo}.

